 
\section{Related Works}\label{sec:related_works}
 
Despite the numerous empirical achievements, unanswered questions concerning the inner workings of neural networks capable of algorithmic reasoning. The ability of self-attention to create low-complexity circuits has been recognized \citep{edelman2022inductive,hahn2020theoretical,merrill2021power}, as well as its capacity to form declarative programs~\citep{weiss2021thinking}, and Turing machines~\citep{dehghani2018universal,giannou2023looped,perez2021attention}. Moreover, it has been demonstrated that interpretable symbolic computations can be drawn from trained models~\citep{clark2019does,tenney2019bert,vig2019visualizing,wang2022interpretability}.

\citet{liu2022transformers} is a closely related work to ours, which studies the expressiveness of low-depth transformers for semi-automata. Their setting corresponds to using only 1 step of CoT and our contribution is to show that allowing more steps of CoT enables the transformers to solve more difficult problems than semi-automata, especially those inherently serial problems, like the circuit value problem, which is $\P$-complete.

\paragraph{Constant precision versus logarithmic precision:} We note that most previous literature on the expressiveness of transformers focuses on the setting of logarithmic precision, including \citep{merrill2023parallelism,merrill2022saturated,merrill2021power,liu2022transformers}, \emph{etc}. One main reason as argued by \citet{merrill2023logic} is that log precision allows the transformer to use uniform attention over the rest tokens. However, recent advancements in LLMs showed that uniform attention might not be necessary towards good performance, at least for natural language tasks. For example, one of the most successful open-sourced LLM, LLAMA2 ~\citep{touvron2023llama} takes the input of a sequence of $4096$ tokens and uses \textsf{BF16} precision, which has 1 sign bit, 8 exponent bits and 7 mantissa bits (plus one extra leading bit). As a consequence, for example, $\mathsf{BF16}$ cannot express any floating-point number between $2^8=256 $ and $2^8+2 = 258$, which makes LLAMA2 impossible to compute uniform attention over $257$ elements.

A concurrent work \citet{feng2023towards} also studies the benefit of CoT via the perspective of expressiveness, where they show with CoT, transformers can solve some specific $\P$-complete problem. Our result is stronger in the sense that we give a simple and clean construction for each problem in $\Ppoly$. We also note the slight difference in the settings, while we mainly focus on constant-precision transformers with $O(\log n)$ embedding size, they focus on $O(\log (n))$ precision transformers with bounded embedding size.
